\documentclass[11pt,a4paper]{article}
\usepackage[margin=25mm]{geometry}
\usepackage{amsmath,fontspec,unicode-math,array}
\setmainfont{Nimbus Roman}
\setmonofont{Nimbus Mono PS}[Scale=MatchLowercase]
\setmathfont{MTPro2Math.otf}[Path=../out/]
% Use upright letters from the active math font for named operators.
\setoperatorfont{\symup}
% Named versions are used only to print alternatives on successive rows.
% A user's document normally selects its features once with setmathfont.
\setmathfont{MTPro2Math.otf}[Path=../out/,version=heavy,
  RawFeature={+cv33}]
\setmathfont{MTPro2Math.otf}[Path=../out/,version=holey,
  RawFeature={+cv34}]
\setmathfont{MTPro2Math.otf}[Path=../out/,version=heavyholey,
  RawFeature={+cv33,+cv34}]
\setmathfont{MTPro2Math.otf}[Path=../out/,version=scriptalt,
  RawFeature={+cv31}]
\setmathfont{MTPro2Math.otf}[Path=../out/,version=curly,
  RawFeature={+ss01}]
\setmathfont{MTPro2Math.otf}[Path=../out/,version=curlyalt,
  RawFeature={+ss01,+cv31}]
\setmathfont{MTPro2Math.otf}[Path=../out/,version=curlytwo,
  RawFeature={+ss02}]
\setmathfont{MTPro2Math.otf}[Path=../out/,version=frakturalt,
  RawFeature={+cv32}]
\setmathfont{MTPro2Math.otf}[Path=../out/,version=straight,
  RawFeature={+cv21}]
\setmathfont{MTPro2Math.otf}[Path=../out/,version=largeparen,
  RawFeature={+cv22}]
\setmathfont{MTPro2Math.otf}[Path=../out/,version=largecurly,
  RawFeature={+cv23}]
\setmathfont{MTPro2Math.otf}[Path=../out/,version=swash,
  RawFeature={+cv01}]
\setmathfont{MTPro2Math.otf}[Path=../out/,version=greekalt,
  RawFeature={+cv02}]
\setmathfont{MTPro2Math.otf}[Path=../out/,version=planckalt,
  RawFeature={+cv03}]
\setmathfont{MTPro2Math.otf}[Path=../out/,version=slantedops,
  RawFeature={+cv11}]
\setmathfont{MTPro2Math.otf}[Path=../out/,version=andalt,
  RawFeature={+cv12}]
\setmathfont{MTPro2Math.otf}[Path=../out/,version=plainhbrace,
  RawFeature={cv24=1}]
\setmathfont{MTPro2Math.otf}[Path=../out/,version=curlyhbrace,
  RawFeature={cv24=2}]
\setmathfont{MTPro2Math.otf}[Path=../out/,version=subsetalt,
  RawFeature={+cv41}]
\setmathfont{MTPro2Math.otf}[Path=../out/,version=orderalt,
  RawFeature={+cv42}]
\setmathfont{MTPro2Math.otf}[Path=../out/,version=shortrel,
  RawFeature={+cv43}]
\setmathfont{MTPro2Math.otf}[Path=../out/,version=thickrel,
  RawFeature={+cv44}]
\setmathfont{MTPro2Math.otf}[Path=../out/,version=curverel,
  RawFeature={+cv45}]
\setmathfont{MTPro2Math.otf}[Path=../out/,version=negone,
  RawFeature={cv46=1}]
\setmathfont{MTPro2Math.otf}[Path=../out/,version=negtwo,
  RawFeature={cv46=2}]
\setmathfont{MTPro2Math.otf}[Path=../out/,version=nsimalt,
  RawFeature={+cv47}]
\setmathfont{MTPro2Math.otf}[Path=../out/,version=emptyalt,
  RawFeature={+cv51}]
\setmathfont{MTPro2Math.otf}[Path=../out/,version=cardsalt,
  RawFeature={+cv52}]
\setmathfont{MTPro2Math.otf}[Path=../out/,version=marksalt,
  RawFeature={+cv53}]
\usepackage[hidelinks,pdftitle={MTPro2 Math: Font guide},
  pdfauthor={}]{hyperref}
\setlength{\parindent}{0pt}
\setlength{\parskip}{5pt}
\setlength{\tabcolsep}{7pt}
\renewcommand{\arraystretch}{1.28}
\emergencystretch=2em
\clubpenalty=10000
\widowpenalty=10000
\newcommand{\code}[1]{\texttt{\detokenize{#1}}}
\newcommand{\sample}[2]{\begingroup\mathversion{#1}$#2$\endgroup}
\newcommand{\bbline}{\symbb{ABGMXYZ\;abcxyz\;0123456789}}
\newcommand{\bbgreek}{\BbbGamma\;\BbbPi\;\Bbbgamma\;\Bbbpi}
\newcommand{\bbitalic}{\mitBbbD\;\mitBbbd\;\mitBbbe\;\mitBbbi\;\mitBbbj}
\newcommand{\scriptline}{\symscr{CGILQSYZrz}}
\newcommand{\boldscriptline}{\symbfscr{CGILQSYZrz}}
\newcommand{\matrixsample}{\left(\begin{matrix}
  2&-1&0&0\\
  -1&2&-1&0\\
  0&-1&2&-1\\
  0&0&-1&2
\end{matrix}\right)}
\newcommand{\tentfunction}{\Lambda(x)=\left\{\begin{aligned}
  1+x,&\quad -1\leq x<0,\\
  1-x,&\quad 0\leq x\leq1,\\
  0,&\quad \text{otherwise}.
\end{aligned}\right.}
% Helpers below arrange specimens; they are not a font interface.
% U+00AE has no named unicode-math command; select the math family explicitly.
\newcommand{\registeredmark}{\Umathchar 0 \symoperators "00AE\relax}
\newcommand{\smallbrace}{\left\{\frac{1}{n}:n\in\symbb{N}\right\}}
\newcommand{\bracepair}{\overbrace{a+b+c}^{s}\quad\underbrace{x+y+z}_{t}}
\newcommand{\operatorsample}{\displaystyle
  \sum_{k=1}^{n}a_k\quad\prod_{k=1}^{n}b_k\quad\coprod_{i\in I}X_i}
\newcommand{\andsample}{p\wedge q\quad\displaystyle\bigwedge_{i=1}^{n}P_i}
\newcommand{\negatedsample}{a\nleq b\quad a\ngeq b\quad
  A\nsubseteq B\quad A\nsupseteq B}
\begin{document}

{\LARGE MTPro2 Math}\par
{\large Font guide}\hfill v1.4\par
\medskip
Use the generated font with LuaLaTeX and \code{unicode-math}:
\begin{verbatim}
\usepackage{unicode-math}
\setmathfont{MTPro2Math.otf}
\end{verbatim}
Use \code{MTPro2MathLite.otf} for Lite. The examples in this guide use Full.
No \code{mtpro2.sty} is needed. Building the font and using it in Word
are described in \code{mtp2otf.pdf}.

\section{Letters and Greek}
Ordinary letters are mathematical italic. Use \code{\symup},
\code{\symit}, \code{\symbfup} and \code{\symbfit} to select a
mathematical alphabet explicitly:
\begin{center}
\begin{tabular}{@{}ll@{}}
\code{\symup{ABGMXYZabcxyz}} & $\symup{ABGMXYZabcxyz}$\\
\code{\symit{ABGMXYZabcxyz}} & $\symit{ABGMXYZabcxyz}$\\
\code{\symbfup{ABGMXYZabcxyz}} & $\symbfup{ABGMXYZabcxyz}$\\
\code{\symbfit{ABGMXYZabcxyz}} & $\symbfit{ABGMXYZabcxyz}$\\
\end{tabular}
\end{center}
The \code{\sym...} commands use the math font rather than the document's
text font. Bold italic requires Full. For Greek, the same selectors give
$\symup{\alpha\beta\Gamma\Omega}$,
$\symit{\alpha\beta\Gamma\Omega}$,
$\symbfup{\alpha\beta\Gamma\Omega}$ and
$\symbfit{\alpha\beta\Gamma\Omega}$.

Character variants select alternative designs for the whole document.
For example, swash z is selected in the preamble:
\begin{verbatim}
\setmathfont{MTPro2Math.otf}[RawFeature={+cv01}]
\end{verbatim}
\begin{center}
\begin{tabular}{@{}lll@{}}
Default & \sample{normal}{z\quad\symbfit{z}} & swash z\\
\code{cv01} & \sample{swash}{z\quad\symbfit{z}} & \\[5pt]
Default & \sample{normal}{\symup{\beta\delta}\quad\symit{\beta\delta}\quad
  \symbfup{\beta\delta}\quad\symbfit{\beta\delta}} & beta, delta\\
\code{cv02} & \sample{greekalt}{\symup{\beta\delta}\quad\symit{\beta\delta}\quad
  \symbfup{\beta\delta}\quad\symbfit{\beta\delta}} & \\[5pt]
Default & \sample{normal}{\hbar} & Planck sign\\
\code{cv03} & \sample{planckalt}{\hbar} & \\
\end{tabular}
\end{center}
These choices are available in both editions; their bold counterparts
require Full. Under \code{unicode-math}, \code{\hbar} and
\code{\hslash} name the same character, whose design is selected by
\code{cv03}.

Features can be combined: \code{RawFeature={+cv01,+cv02}} selects both.
For a feature with numbered choices, use, for example,
\code{RawFeature={cv24=2}}. Omit a feature to keep its default design.

\newpage
\section{Blackboard}
Blackboard is available in Full. Use \code{\symbb{...}} for double-struck letters and digits.
The default design and three upright alternatives are shown below.

\begin{center}
\code{\symbb{ABGMXYZ\;abcxyz\;0123456789}}\par\smallskip
\begin{tabular}{@{}ll@{}}
Default & \sample{normal}{\bbline}\\
\code{cv33} & \sample{heavy}{\bbline}\\
\code{cv34} & \sample{holey}{\bbline}\\
\code{cv33,cv34} & \sample{heavyholey}{\bbline}\\
\end{tabular}
\end{center}
\code{cv33} selects heavier upright letters. \code{cv34} selects the
Holey Roman family; combine them for its heavier upright design.
For example:
\begin{verbatim}
\setmathfont{MTPro2Math.otf}[RawFeature={+cv33,+cv34}]
\end{verbatim}

Four Greek characters and five italic characters have standard Unicode
inputs. \code{cv34} changes these to the corresponding Holey designs;
\code{cv33} leaves them unchanged.

\begin{center}
\begin{tabular}{@{}lll@{}}
 & Greek & Italic\\
Default & \sample{normal}{\bbgreek} & \sample{normal}{\bbitalic}\\
\code{cv34} & \sample{holey}{\bbgreek} & \sample{holey}{\bbitalic}\\
\end{tabular}
\end{center}
\begin{verbatim}
\BbbGamma \BbbPi \Bbbgamma \Bbbpi
\mitBbbD \mitBbbd \mitBbbe \mitBbbi \mitBbbj
\end{verbatim}
The italic Blackboard selection is limited to these five characters.

\newpage
\section{Script and Fraktur}
These alphabets are available in Full. Use \code{\symscr} and \code{\symbfscr} for regular and bold Script.
\code{cv31} changes selected letters in both weights. Curly Script is
selected with \code{ss01}; it also has alternatives under \code{cv31}.
\code{ss02} selects a second Curly series directly.

\begin{center}
\code{\symscr{CGILQSYZrz}}\par\smallskip
\begin{tabular}{@{}ll@{}}
Default & \sample{normal}{\scriptline}\\
\code{cv31} & \sample{scriptalt}{\scriptline}\\
\code{ss01} & \sample{curly}{\scriptline}\\
\code{ss01,cv31} & \sample{curlyalt}{\scriptline}\\
\code{ss02} & \sample{curlytwo}{\scriptline}\\
\end{tabular}
\end{center}
\begin{verbatim}
\setmathfont{MTPro2Math.otf}[RawFeature={+ss01,+cv31}]
\end{verbatim}
\code{ss02} takes precedence over the other regular Script choices.
The Curly selectors affect regular Script; bold Script retains its own
design, with its alternatives controlled by \code{cv31}.

\begin{center}
\code{\symbfscr{CGILQSYZrz}}\par\smallskip
\begin{tabular}{@{}ll@{}}
Default & \sample{normal}{\boldscriptline}\\
\code{cv31} & \sample{scriptalt}{\boldscriptline}\\
\end{tabular}
\end{center}

Fraktur is selected with \code{\symfrak} or \code{\symbffrak}.
Its alternatives use \code{cv32}, independently of the Script choices.

\begin{center}
\begin{tabular}{@{}lll@{}}
 & \code{\symfrak{IRYxy}} & \code{\symbffrak{IRYxy}}\\
Default & \sample{normal}{\symfrak{IRYxy}}
  & \sample{normal}{\symbffrak{IRYxy}}\\
\code{cv32} & \sample{frakturalt}{\symfrak{IRYxy}}
  & \sample{frakturalt}{\symbffrak{IRYxy}}\\
\end{tabular}
\end{center}
\begin{verbatim}
\setmathfont{MTPro2Math.otf}[RawFeature={+cv31,+cv32}]
\end{verbatim}

\newpage
\section{Parentheses and braces}
Use \code{\left} and \code{\right} to size delimiters to a formula.
The default series uses the original ordinary growth range, followed by
an assembly of repeatable parts. The default braces use the morphed
design, which approaches the assembled shape as it grows.

\code{cv22} extends parentheses through the large ready-made designs.
The two examples below use the same input and font size.

\begin{center}
\begin{tabular}{@{}ll@{}}
Default & \sample{normal}{\matrixsample}\\[10pt]
\code{cv22} & \sample{largeparen}{\matrixsample}\\
\end{tabular}
\end{center}
\begin{verbatim}
\setmathfont{MTPro2Math.otf}[RawFeature={+cv22}]
\end{verbatim}

For braces, \code{cv21} selects straight braces with the ordinary growth
range. \code{cv23} selects Curly braces with an extended ready-made range.

At larger sizes the default and straight series use the same assembly.
Each row shows a small set and a larger piecewise function:
\begin{center}
\begin{tabular}{@{}lll@{}}
Default & \sample{normal}{\smallbrace}
  & \sample{normal}{\tentfunction}\\[10pt]
\code{cv21} & \sample{straight}{\smallbrace}
  & \sample{straight}{\tentfunction}\\[10pt]
\code{cv23} & \sample{largecurly}{\smallbrace}
  & \sample{largecurly}{\tentfunction}\\
\end{tabular}
\end{center}
\begin{verbatim}
\setmathfont{MTPro2Math.otf}[RawFeature={+cv22,+cv23}]
\end{verbatim}
Use paired delimiters, or end with \code{\right.} for a one-sided brace:
\begin{verbatim}
\left( ... \right)     \left\{ ... \right\}
\left\{ ... \right.
\end{verbatim}

The large series reaches 28.8\,em before continuing with the ordinary
assembly. \code{cv21} takes precedence over \code{cv23}; choose one brace
style. Parenthesis selection remains independent.


\newpage
\section{Accents and horizontal braces}
Use a fixed accent for a letter and a wide accent for a group. These
constructions are available in both editions.
\begin{center}
\begin{tabular}{@{}ll@{}}
\code{\hat{x}, \widehat{ABC}} & $\hat{x}\quad\widehat{ABC}$\\[2pt]
\code{\tilde{x}, \widetilde{ABC}} & $\tilde{x}\quad\widetilde{ABC}$\\[2pt]
\code{\check{x}, \widecheck{ABC}} & $\check{x}\quad\widecheck{ABC}$\\[2pt]
\code{\bar{M}, \wideoverbar{M+N}} & $\bar{M}\quad\wideoverbar{M+N}$\\[2pt]
\code{\vec{v}, \overrightarrow{AB}} & $\vec{v}\quad\overrightarrow{AB}$\\[2pt]
\code{\overleftarrow{AB}} & $\overleftarrow{AB}$\\[2pt]
\code{\dot{x}, \ddot{x}, \hat{i}} & $\dot{x}\quad\ddot{x}\quad\hat{i}$\\
\end{tabular}
\end{center}
\code{\bar} and \code{\overbar} select the fixed accent;
\code{\wideoverbar} selects the stretching accent.
\code{\overline} draws a rule instead. An operator such as
$\hat{H}$ uses the ordinary \code{\hat{H}}.

Horizontal parentheses and braces cover their contents. For example:
\begin{verbatim}
\overparen{AB} \qquad \underbrace{V\times\cdots\times V}_{n}
\end{verbatim}
\[
  \overparen{AB}\qquad\underbrace{V\times\cdots\times V}_{n}
\]
\code{cv24} selects the design of both overbraces and underbraces:
\begin{center}
\begin{tabular}{@{}ll@{}}
Default & \sample{normal}{\bracepair}\\[10pt]
\code{cv24=1}: plain & \sample{plainhbrace}{\bracepair}\\[10pt]
\code{cv24=2}: Curly & \sample{curlyhbrace}{\bracepair}\\
\end{tabular}
\end{center}
\begin{verbatim}
\setmathfont{MTPro2Math.otf}[RawFeature={cv24=2}]
\end{verbatim}
This choice is independent of the vertical-brace features. Use:
\begin{verbatim}
\overbrace{...}^{...}     \underbrace{...}_{...}
\end{verbatim}

\newpage
\section{Operators and formula sizes}
To use the math font's upright letters for named operators such as
\code{\sin} and \code{\cos}, set:
\begin{verbatim}
\setoperatorfont{\symup}
\end{verbatim}
Large operators have text and display designs. The math style selects
them automatically; \code{\limits} and \code{\nolimits} control the
placement of limits.
\[
  \textstyle\sum_{k=1}^{n}a_k\quad\int_0^1 f(x)\,dx
  \qquad\displaystyle\sum_{k=1}^{n}a_k\quad\int_0^1 f(x)\,dx
\]
\code{cv11} selects slanted sums, products and coproducts in both sizes.
\begin{center}
\begin{tabular}{@{}ll@{}}
Default & \sample{normal}{\operatorsample}\\[9pt]
\code{cv11} & \sample{slantedops}{\operatorsample}\\
\end{tabular}
\end{center}
\begin{verbatim}
\setmathfont{MTPro2Math.otf}[RawFeature={+cv11}]
\end{verbatim}
\code{cv12} changes conjunction signs, both \code{\wedge} and
\code{\bigwedge}:
\begin{center}
\begin{tabular}{@{}ll@{}}
Default & \sample{normal}{\andsample}\\[9pt]
\code{cv12} & \sample{andalt}{\andsample}\\
\end{tabular}
\end{center}
Both operator features are available in Full and Lite.
Other operators retain their own designs:
\[
  \bigcap_{i=1}^{n}A_i\qquad\bigcup_{i=1}^{n}A_i
  \qquad\bigotimes_{i=1}^{n}V_i\qquad\oint_C\vec{F}\cdot d\vec{r}
\]
Subscripts and superscripts select the source's optical sizes without
an explicit feature setting:
\[
  x+\alpha\qquad 2^{x+\alpha}\qquad 2^{2^{x+\alpha}}
\]
Radicals and fractions use the ordinary commands:
\begin{verbatim}
\sqrt{\frac{a^2+b^2}{c^2+d^2}} \qquad \sqrt[3]{x+y}
\end{verbatim}
\[
  \sqrt{\frac{a^2+b^2}{c^2+d^2}}\qquad\sqrt[3]{x+y}
\]
The original package's explicit extra-large operator commands are not
provided.

\newpage
\section{Relations}
The following Full-edition features change relation designs without
changing their input commands or mathematical meanings.
\begin{center}
\begin{tabular}{@{}lll@{}}
Default & \sample{normal}{\subsetneq\;\supsetneq\;\subsetneqq\;\supsetneqq} & \\
\code{cv41} & \sample{subsetalt}{\subsetneq\;\supsetneq\;\subsetneqq\;\supsetneqq}
  & \code{\subsetneq \supsetneq}\\
 & & \code{\subsetneqq \supsetneqq}\\[5pt]
Default & \sample{normal}{\lneqq\;\gneqq} & \\
\code{cv42} & \sample{orderalt}{\lneqq\;\gneqq}
  & \code{\lneqq \gneqq}\\[5pt]
Default & \sample{normal}{\mid\;\nmid\;\parallel\;\nparallel} & \\
\code{cv43} & \sample{shortrel}{\mid\;\nmid\;\parallel\;\nparallel}
  & \code{\mid \nmid \parallel \nparallel}\\[5pt]
Default & \sample{normal}{\sim\;\approx} & \\
\code{cv44} & \sample{thickrel}{\sim\;\approx}
  & \code{\sim \approx}\\[5pt]
Default & \sample{normal}{\propto\;\smile\;\frown} & \\
\code{cv45} & \sample{curverel}{\propto\;\smile\;\frown}
  & \code{\propto \smile \frown}\\
\end{tabular}
\end{center}
\code{cv43} selects shorter relation bars; it does not change the
delimiter bars in \code{\left|x\right|} or
\code{\left\Vert v\right\Vert}. \code{cv44} selects thicker
similarity and approximation signs.

Negated comparisons and inclusions have two further choices under
\code{cv46}. Use their direct commands:
\begin{verbatim}
a\nleq b \quad a\ngeq b \quad A\nsubseteq B \quad A\nsupseteq B
\end{verbatim}
\begin{center}
\begin{tabular}{@{}ll@{}}
Default & \sample{normal}{\negatedsample}\\[4pt]
\code{cv46=1} & \sample{negone}{\negatedsample}\\[4pt]
\code{cv46=2} & \sample{negtwo}{\negatedsample}\\
\end{tabular}
\end{center}
The not-similar sign is independent:
\begin{center}
\begin{tabular}{@{}lll@{}}
Default & \sample{normal}{a\nsim b} & \code{a\nsim b}\\
\code{cv47} & \sample{nsimalt}{a\nsim b} & \\
\end{tabular}
\end{center}
\begin{verbatim}
\setmathfont{MTPro2Math.otf}[RawFeature={cv46=1,+cv47}]
\end{verbatim}

\newpage
\section{Other symbols}
\code{cv51} selects the alternative empty-set sign in Full.
\code{\emptyset} and \code{\varnothing} refer to the same Unicode
character under \code{unicode-math}.
\begin{center}
\begin{tabular}{@{}lll@{}}
Default & \sample{normal}{\emptyset} & \code{\emptyset}\\
\code{cv51} & \sample{emptyalt}{\emptyset} & \\[5pt]
Default & \sample{normal}{\spadesuit\;\clubsuit}
  & \code{\spadesuit \clubsuit}\\
\code{cv52} & \sample{cardsalt}{\spadesuit\;\clubsuit} & \\[5pt]
Default & \sample{normal}{\mathyen\;\registeredmark}
  & \code{\mathyen}, U+00AE\\
\code{cv53} & \sample{marksalt}{\mathyen\;\registeredmark} & \\
\end{tabular}
\end{center}
\code{cv52} gives shaded spades and clubs in both editions.
\code{cv53} changes the yen and registered signs in Full.

The usual symbol commands include \code{\ell}, \code{\partial},
\code{\nabla}, \code{\infty}, \code{\dagger} and \code{\ddagger}:
\[
  \ell\qquad\partial\qquad\nabla\qquad\infty\qquad\dagger\qquad\ddagger
\]
Mathematical sans-serif and monospace alphabets, a complete bold-symbol
version, and a compatibility package for MTPro2-specific macros are
not included.

\newpage
\section{Three math sizes}
Each group repeats an expression at normal size, in a superscript and in
an exponent of a superscript. The exponent construction is used here
only to compare the three optical sizes.

\[
  A,\ldots,Z\quad a,\ldots,z\quad\Gamma,\Omega\quad\alpha_i,\omega^j
\]
\[
  2^{A,\ldots,Z\quad a,\ldots,z\quad\Gamma,\Omega\quad\alpha_i,\omega^j}
\]
\[
  2^{2^{A,\ldots,Z\quad a,\ldots,z\quad\Gamma,\Omega\quad\alpha_i,\omega^j}}
\]

\[
  (f\circ g)'(x)=f'(g(x))\,g'(x)
\]
\[
  2^{(f\circ g)'(x)=f'(g(x))\,g'(x)}
\]
\[
  2^{2^{(f\circ g)'(x)=f'(g(x))\,g'(x)}}
\]

\[
  \sum_{k=1}^{n}\frac{1}{k(k+1)}=\frac{n}{n+1}
\]
\[
  2^{\sum_{k=1}^{n}\frac{1}{k(k+1)}=\frac{n}{n+1}}
\]
\[
  2^{2^{\sum_{k=1}^{n}\frac{1}{k(k+1)}=\frac{n}{n+1}}}
\]

\[
  \int_{-\infty}^{\infty}e^{-x^2}\,dx=\sqrt{\pi}
\]
\[
  2^{\int_{-\infty}^{\infty}e^{-x^2}\,dx=\sqrt{\pi}}
\]
\[
  2^{2^{\int_{-\infty}^{\infty}e^{-x^2}\,dx=\sqrt{\pi}}}
\]

\[
  \omega=v\,dx+w\,dy,\qquad d\omega=\left(\frac{\partial w}{\partial x}-\frac{\partial v}{\partial y}\right)\,dx\wedge dy
\]
\[
  2^{\omega=v\,dx+w\,dy,\qquad d\omega=\left(\frac{\partial w}{\partial x}-\frac{\partial v}{\partial y}\right)\,dx\wedge dy}
\]
\[
  2^{2^{\omega=v\,dx+w\,dy,\qquad d\omega=\left(\frac{\partial w}{\partial x}-\frac{\partial v}{\partial y}\right)\,dx\wedge dy}}
\]

\[
  R_{ijkl}=-R_{jikl}=-R_{ijlk}=R_{klij}
\]
\[
  2^{R_{ijkl}=-R_{jikl}=-R_{ijlk}=R_{klij}}
\]
\[
  2^{2^{R_{ijkl}=-R_{jikl}=-R_{ijlk}=R_{klij}}}
\]

\end{document}
